% !PRINTED 633
\weberpage{633}
Diese Resolventen sind ganze algebraische Zahlen in gewissen
Kreistheilungskörpern, und wir haben, wenn $\varepsilon^\lambda$ nicht $=1$ ist,
\[
(\varepsilon^\lambda,\varrho)\,(\varepsilon^{-\lambda},\varrho) = (-1)^\lambda p,
\]
und für $\lambda = m'n$
\weq{5}{(r^n,\eta)\,(r^{-n},\eta) = \pm\, p.}

Es sind also alle diese Resolventen Theiler der natürlichen
Primzahl $p$.

Ausser diesen Resolventen haben wir an der angeführten
Stelle noch andere ganze Zahlen der Kreistheilungskörper abgeleitet, nämlich
\weq{6}{\psi_{\lambda,\mu}(\varepsilon) = \sum\limits_{1,\,p-2}^{t} \varepsilon^{\mu\,\mathrm{ind}\,t - (\lambda+\mu)\,\mathrm{ind}\,(t+1)},}
(worin sich die Indices auf die Basis $g$ beziehen), die im Allgemeinen dem Körper $\Omega_{p-1}$ angehören, in dem besonderen Falle, wo $\lambda$ und $\mu$ durch $m'$ theilbar sind, dem Körper $\Omega_m$. Von besonderer Wichtigkeit sind darunter die Zahlen des Körpers $\Omega_m$:
\weq{7}{\psi_{sm',\,m'}(\varepsilon) = \psi_s(r) = \sum\limits_{1,\,p-2}^{t} r^{\mathrm{ind}\,t - (s+1)\,\mathrm{ind}\,(t+1)}.}

Auch diese Zahlen $\psi_{\lambda,\mu}(\varepsilon)$ sind, wenn keine der Zahlen $\lambda$, $\mu$, $\lambda+\mu$ durch $p-1$ theilbar ist, Theiler von $p$ nach der Formel
\[
\psi_{\lambda,\mu}(\varepsilon)\,\psi_{\lambda,\mu}(\varepsilon^{-1}) = p,
\]
aus der sich wieder nach (7) ergiebt:
\weq{8}{\psi_s(r)\,\psi_s(r^{-1}) = p,}
vorausgesetzt, dass weder $s$ noch $s+1$ durch $m$ theilbar ist. Gewisse Verbindungen der Resolventen (4), darunter die $m^{\mathrm{te}}$ Potenz, haben wir durch die Functionen $\psi$ ausgedrückt [Bd.\ I, \S.\,169, (14), (15)]:
\weq{9}{(r,\eta)^n\,(r^n,\eta)^{-1} = \psi_1(r)\,\psi_2(r)\cdots\psi_{n-1}(r),}
so lange $n<m$ ist, und
\weq{10}{(r,\eta)\,(r^{-1},\eta) = \pm\,p}
\weq{11}{(r,\eta)^m = (-1)^{\frac{p-1}{m}}\, p\, \psi_1(r)\,\psi_2(r)\cdots\psi_{m-2}(r),}
woraus hervorgeht, dass diese $m^{\mathrm{ten}}$ Potenzen dem Körper $\Omega_m$ angehören.

Wir bezeichnen nun, wie früher, mit $\mathfrak{p}_n$ die verschiedenen Primtheiler von $p$ im Körper $\Omega_m$. Die Zahlen $\psi_s(r)$, die gleichfalls diesem Körper angehören und Factoren von $p$ sind, müssen
%CONT>>
